% !TEX program = xelatex
% 上海电机学院 本科生毕业设计（论文）答辩记录表 示例（文字为中性占位）
% 编译：latexmk defense.tex
% 开题答辩、中期检查答辩、毕业设计（论文）答辩各填一份（2026 版要素 3.1 归档目录 (9)(10)(11)），
% 三份表式相同，只改 \defensetype。答辩记录由答辩组成员填写，不得由学生填写（3.3.5(4)）。
\documentclass{sdjudefense}

% ---------------- 表头 ----------------
% 表中每个栏目都在源码中直接填写；没有内容的留空（花括号内不写），版面上照原表留出空白。
\title{兆瓦级风力发电机控制系统设计——液压系统}
\author{XXX}
\studentid{XXXXXXXXX}
\major{电气工程及其自动化}
\classname{XXXXXXX}
\defensetype{proposal}           % 答辩类型（勾选）：proposal 开题答辩 / midterm 中期检查答辩 / thesis 毕业设计（论文）答辩，也可写中文
\committee{XXX（组长）、XXX、XXX}  % 专家组成员
\recorder{XXX}                   % 记录人；留空则只印「记录人：」（签名手签，或由系统放入签名图片）
\recorddate{20XX 年 XX 月 XX 日}   % 日期；留空为空白的「年 月 日」

\begin{document}

\maketitle

% 一问一答，只记问题，不记评价；答辩人的回答详细记录（3.3.5(3)）
\section{答辩记录}
% 每个问题一行 \question{问题}，紧接着 \answer{回答}
\question{液压系统的工作压力是按什么工况确定的？}
\answer{按变桨油缸在额定风速下所需的最大推力，结合油缸有效面积计算，再留出一定裕量……}

\question{比例阀的选型依据是什么？}
\answer{依据系统所需的最大流量和响应时间，对比了两种型号的额定流量与频响……}

\question{下一阶段如何验证控制程序？}
\answer{先在仿真平台上进行变桨距控制的闭环仿真，再在实验台上调试……}

\end{document}
